\begin{derivation}[从\eqnrefs{eq:one-param-group-dp11}到\eqnrefs{eq:matrix-lie-bracket-math-dp76}]
先沿固定切向方向 $X=\dot g(0)$ 考察单参数子群。由群关系 $g(t+s)=g(t)g(s)$ 对 $s$ 求导并令 $s=0$，得到
\begin{align*}
 \dot g(t)
 &=g(t)\left.\frac{\dd g(s)}{\dd s}\right|_{s=0}\\
 &=g(t)X.
\end{align*}
交换 $t,s$ 的角色，同理由 $g(t+s)=g(s)g(t)$ 得 $\dot g(t)=Xg(t)$，所以 $X$ 与该单参数子群上的 $g(t)$ 对易。设
\begin{equation*}
 g(t)=\sum_{n=0}^{\infty}t^nG_n,
\end{equation*}
代入 $\dot g=Xg$ 并比较 $t^n$ 的系数，得到
\begin{equation*}
 (n+1)G_{n+1}=XG_n,\qquad G_0=I.
\end{equation*}
递推给出 $G_n=X^n/n!$，于是
\begin{equation*}
 g(t)=I+tX+\frac{t^2X^2}{2!}+\cdots=\ee^{tX},
\end{equation*}
即正文\eqnrefs{eq:one-param-exp-dp11}。

接着取另一个切向量 $Y\in\mathfrak g=T_eG$，并选一条群曲线 $h(t)$ 满足 $h(0)=I$、$\dot h(0)=Y$。由刚得到的 $\ee^{sX}\in G$ 以及群乘法和取逆的封闭性，
\begin{equation*}
 h_s(t)=\ee^{sX}h(t)\ee^{-sX}\in G.
\end{equation*}
因此它在 $t=0$ 的切向量仍属于单位元切空间：
\begin{equation*}
 Y(s)=\left.\frac{\dd h_s(t)}{\dd t}\right|_{t=0}
 =\ee^{sX}Ye^{-sX}\in\mathfrak g.
\end{equation*}
对 $s$ 再求导，
\begin{align*}
 \frac{\dd Y(s)}{\dd s}
 &=Xe^{sX}Ye^{-sX}-\ee^{sX}Ye^{-sX}X,\\
 \left.\frac{\dd Y(s)}{\dd s}\right|_{s=0}
 &=XY-YX.
\end{align*}
因为 $Y(s)$ 是向量空间 $\mathfrak g$ 内的光滑曲线，其导数仍属于 $\mathfrak g$，所以
\begin{equation*}
 XY-YX\in\mathfrak g,
\end{equation*}
从而得到正文\eqnrefs{eq:matrix-lie-bracket-math-dp76}。
\end{derivation}

\begin{derivation}[从\eqnrefs{eq:lie-algebra-general-dp11}到\eqnrefs{eq:rep-lie-algebra-dp11}]
需要验证的关键点是：群表示在单位元处求导以后，仍然保持 Lie 括号。取两个抽象生成方向 $X,Y\in\mathfrak g$。群中的共轭变换满足
\[
\operatorname{Ad}_{\ee^{sX}}Y=\ee^{sX}Ye^{-sX}.
\]
表示保持群乘法，因此
\begin{align*}
D_R(\ee^{sX})D_R(\ee^{tY})D_R(\ee^{-sX})
&=D_R\!\left(\ee^{sX}\ee^{tY}\ee^{-sX}\right).
\end{align*}
先对 $t$ 在 $t=0$ 求导，得到
\begin{align*}
D_R(\ee^{sX})\,dD_R(Y)\,D_R(\ee^{-sX})
&=dD_R\!\left(\operatorname{Ad}_{\ee^{sX}}Y\right).
\end{align*}
再对 $s$ 在 $s=0$ 求导。左边给矩阵对易子，右边给抽象伴随作用的导数：
\begin{align*}
[dD_R(X),dD_R(Y)]
&=dD_R([X,Y]_{\mathfrak g}).
\end{align*}
现在令 $X=X_a$、$Y=X_b$，并代入正文\eqnrefs{eq:lie-algebra-general-dp11}：
\begin{align*}
[dD_R(X_a),dD_R(X_b)]
&=f_{ab}{}^c\,dD_R(X_c).
\end{align*}
对酉表示，正文已定义 $dD_R(X_a)=-\ii T_R^a$，于是
\begin{align*}
[-\ii T_R^a,-\ii T_R^b]
&=f_{ab}{}^c(-\ii T_R^c),\\
-[T_R^a,T_R^b]
&=-\ii f_{ab}{}^cT_R^c.
\end{align*}
两边乘以 $-1$，得到正文\eqnrefs{eq:rep-lie-algebra-dp11}：
\[
[T_R^a,T_R^b]=\ii f_{ab}{}^cT_R^c.
\]
\end{derivation}

% DeepPass21: detailed Lie-algebra derivations gathered at the end of the owner subsection.










\begin{derivation}[从\eqnrefs{eq:generator-trace-normalization-dp11}到\eqnrefs{eq:suN-fierz-math-dp68}]



正文已经固定基本表示生成元的无迹性和迹归一化。四指标生成元和只能由两个 Kronecker 张量构成；用两种收缩固定它们的系数，就得到正文\eqnrefs{eq:suN-fierz-math-dp68}。设 $SU(N)$ 生成元满足
\begin{equation*}
 \Tr(T^aT^b)=T_F\delta^{ab},
 \qquad \Tr T^a=0.
\end{equation*}
把所有无迹厄米生成元与单位矩阵一起看成 $N\times N$ 矩阵空间的一组正交基。定义
\begin{equation*}
 X_{ij;kl}\equiv\sum_{a=1}^{N^2-1}(T^a)_{ij}(T^a)_{kl}.
\end{equation*}
由于 $i,l$ 与 $j,k$ 分别形成基本表示--反基本表示指标对，$SU(N)$ 不变量只能由
\begin{equation*}
 \delta_{il}\delta_{jk},
 \qquad
 \delta_{ij}\delta_{kl}
\end{equation*}
组成，所以写
\begin{equation*}
 X_{ij;kl}=A\delta_{il}\delta_{jk}+B\delta_{ij}\delta_{kl}.
\end{equation*}

第一组收缩取 $i=j$ 并求和：
\begin{align*}
 \sum_iX_{ii;kl}
 &=\sum_a\Tr(T^a)(T^a)_{kl}=0,\\
 &=A\sum_i\delta_{il}\delta_{ik}
 +B\sum_i\delta_{ii}\delta_{kl},\\
 &=(A+NB)\delta_{kl}.
\end{align*}
所以
\begin{equation*}
 B=-\frac AN.
\end{equation*}
第二组收缩取 $i=l$、$j=k$ 并对 $i,j$ 求和：
\begin{align*}
 \sum_{ij}X_{ij;ji}
 &=\sum_a\Tr(T^aT^a),\\
 &=T_F(N^2-1).
\end{align*}
而拟设右边给
\begin{align*}
 \sum_{ij}
 \left(A\delta_{ii}\delta_{jj}+B\delta_{ij}\delta_{ji}\right)
 &=AN^2+BN,\\
 &=A(N^2-1).
\end{align*}
故 $A=T_F$，从而
\begin{equation*}
 {
 \sum_a(T^a)_{ij}(T^a)_{kl}
 =T_F\left(
 \delta_{il}\delta_{jk}
 -\frac1N\delta_{ij}\delta_{kl}
 \right).}
\end{equation*}

检查：令 $k=j$ 并求和 $j$，左边是 $(\sum_aT^aT^a)_{il}=C_F\delta_{il}$；右边给
\begin{equation*}
 T_F\left(N-\frac1N\right)\delta_{il}
 =T_F\frac{N^2-1}{N}\delta_{il},
\end{equation*}
所以 $C_F=T_F(N^2-1)/N$，对 $T_F=1/2$ 恢复 $C_F=(N^2-1)/(2N)$。
\end{derivation}

\begin{derivation}[从\eqnrefs{eq:suN-fierz-math-dp68}到\eqnrefs{eq:ca-general-math-dp11}]



伴随 Casimir 可以不写任何显式 $SU(N)$ 矩阵，而由刚得到的完备关系求出。比较同一个对易子迹的两种计算即可固定 $C_A$。从基本对易子
\begin{equation*}
 [T^a,T^b]=\ii f^{abc}T^c
\end{equation*}
出发。考虑
\begin{equation*}
 \Tr([T^a,T^c][T^b,T^c]),
\end{equation*}
并对重复的 $c$ 求和。一方面直接代入结构常数：
\begin{align*}
 \sum_c\Tr([T^a,T^c][T^b,T^c])
 &=\sum_c\Tr(\ii f^{acd}T^d\,\ii f^{bce}T^e),\\
 &=-T_F f^{acd}f^{bcd},\\
 &=-T_F C_A\delta^{ab}.
\end{align*}
另一方面展开对易子：
\begin{align*}
 \sum_c\Tr([T^a,T^c][T^b,T^c])
 ={}&\sum_c\operatorname{Tr}\!\left[
 T^aT^cT^bT^c-T^aT^cT^cT^b\right.\\
 &\left.\hspace{2.5cm}-T^cT^aT^bT^c+T^cT^aT^cT^b
 \right].
\end{align*}
逐项用 Fierz/完备关系。对任意矩阵 $X$，由 Fierz 恒等式可得
\begin{equation*}
 \sum_c T^cXT^c
 =T_F\left[(\Tr X)I-\frac1N X\right].
\end{equation*}
因此
\begin{align*}
 \sum_c\Tr(T^aT^cT^bT^c)
 &=-\frac{T_F}{N}\Tr(T^aT^b),\\
 \sum_c\Tr(T^cT^aT^cT^b)
 &=-\frac{T_F}{N}\Tr(T^aT^b).
\end{align*}
另外 $\sum_cT^cT^c=C_FI$，故中间两项各为
\begin{equation*}
 -C_F\Tr(T^aT^b).
\end{equation*}
合并：
\begin{align*}
 \sum_c\Tr([T^a,T^c][T^b,T^c])
 &=-2\left(C_F+\frac{T_F}{N}\right)
 \Tr(T^aT^b),\\
 &=-2T_F\left(C_F+\frac{T_F}{N}\right)\delta^{ab}.
\end{align*}
取 $T_F=1/2$、$C_F=(N^2-1)/(2N)$：
\begin{equation*}
 2\left(C_F+\frac{T_F}{N}\right)
 =2\left(\frac{N^2-1}{2N}+\frac1{2N}\right)=N.
\end{equation*}
与第一种计算比较即得
\begin{equation*}
 C_A=N.
\end{equation*}
这条推导清楚展示了伴随 Casimir 与基本完备关系之间的关系。
\end{derivation}





















